\begin{proof}[Proof of \cref{obj:prop:tie-face-collapse}]
Fix a cell \(x\) with \(p_x>0\), and let \(c\) be the observable capacity vector from \(P_{\mathrm{obs}}\). By \cref{obj:prop:capacity-identification}, \(c\) is valid, so
\[
\ell_i(x)\ge 0,\qquad h_j(x)\ge 0
\]
for all \(i,j\in\mathcal Y\), and
\[
\Delta q(x)
=
P(S_1=1,C\mid X=x)-P(S_0=1,C\mid X=x).
\]
Also, by \cref{obj:def:observable-capacities},
\[
q_0(x)=\sum_i \ell_i(x),\qquad
q_1(x)=\sum_j h_j(x),\qquad
\Delta q(x)=q_1(x)-q_0(x),\qquad
m(x)=\min\{q_0(x),q_1(x)\}.
\]
For a coupling array \(\gamma_x\), write
\[
\operatorname{row}_i(\gamma_x)=\sum_j \gamma_{ij,x},
\qquad
\operatorname{col}_j(\gamma_x)=\sum_i \gamma_{ij,x},
\qquad
\operatorname{tot}(\gamma_x)=\sum_i\sum_j\gamma_{ij,x}.
\]
These totals satisfy
\[
\sum_i\operatorname{row}_i(\gamma_x)=\operatorname{tot}(\gamma_x),
\qquad
\sum_j\operatorname{col}_j(\gamma_x)=\operatorname{tot}(\gamma_x).
\]
% lean: capacity_identification

\begin{enumerate}
\item We first prove the face equivalence. If \(\gamma_x\in\Gamma_x^\ast\), then \cref{obj:def:partial-transport-polytope} gives
\[
\operatorname{row}_i(\gamma_x)\le \ell_i(x),\qquad
\operatorname{col}_j(\gamma_x)\le h_j(x),\qquad
\operatorname{tot}(\gamma_x)=m(x).
\]
When \(q_0(x)\le q_1(x)\),
\[
\sum_i\{\ell_i(x)-\operatorname{row}_i(\gamma_x)\}
=
q_0(x)-m(x)
=
0.
\]
Each summand is nonnegative, hence
\[
\operatorname{row}_i(\gamma_x)=\ell_i(x)\qquad\text{for every }i.
\]
Similarly, when \(q_1(x)\le q_0(x)\),
\[
\operatorname{col}_j(\gamma_x)=h_j(x)\qquad\text{for every }j.
\]

By the cellwise output of \cref{obj:def:threshold-flow-construction}, choose a feasible exact-mass coupling
\[
\gamma_x^{\mathrm L}\in \Gamma_x^\ast .
\]
Assume
\[
\Gamma_x^{+}=\Gamma_x^{-},
\qquad
\Gamma_x^{-}=\Gamma_x^{0}.
\]
If \(q_0(x)\le q_1(x)\), the preceding row-exactness argument puts
\(\gamma_x^{\mathrm L}\) in \(\Gamma_x^{+}\). The displayed face equalities then put
\(\gamma_x^{\mathrm L}\) in \(\Gamma_x^0\). Therefore
\[
q_0(x)=\operatorname{tot}(\gamma_x^{\mathrm L})=q_1(x),
\]
using exact rows for the first equality and exact columns for the second. If \(q_1(x)\le q_0(x)\), the same argument with columns first puts
\(\gamma_x^{\mathrm L}\) in \(\Gamma_x^{-}\), then in \(\Gamma_x^0\), and again gives
\[
q_0(x)=\operatorname{tot}(\gamma_x^{\mathrm L})=q_1(x).
\]
Thus \(\Delta q(x)=0\).

Conversely, suppose \(\Delta q(x)=0\). Then
\[
q_0(x)=q_1(x)=m(x).
\]
If \(\gamma_x\in\Gamma_x^{+}\), then its exact row constraints give
\[
\operatorname{tot}(\gamma_x)
=
\sum_i\operatorname{row}_i(\gamma_x)
=
\sum_i\ell_i(x)
=
q_0(x)
=
m(x),
\]
so \(\gamma_x\in\Gamma_x^\ast\). Since \(q_1(x)\le q_0(x)\), the column-exactness argument above gives
\[
\operatorname{col}_j(\gamma_x)=h_j(x)\qquad\text{for every }j,
\]
and hence \(\gamma_x\in\Gamma_x^{-}\). The reverse inclusion follows symmetrically from exact columns and row-exactness. Therefore
\[
\Gamma_x^{+}=\Gamma_x^{-}.
\]
The same reasoning sends every element of \(\Gamma_x^{-}\) to \(\Gamma_x^0\), while \(\Gamma_x^0\subseteq\Gamma_x^{-}\) is immediate from the defining constraints. Hence
\[
\Gamma_x^{-}=\Gamma_x^0.
\]
This proves
\[
\bigl(\Gamma_x^{+}=\Gamma_x^{-}\ \text{and}\ \Gamma_x^{-}=\Gamma_x^{0}\bigr)
\Longleftrightarrow
\Delta q(x)=0.
\]
% lean: polytope_tie_iff_gap_zero

\item Now assume \(\Delta q(x)=0\). Since \(p_x>0\), conditional probabilities given \(X=x\) use the denominator \(p_x\). The displayed identity from \cref{obj:prop:capacity-identification} yields
\[
P(S_1=1,C\mid X=x)=P(S_0=1,C\mid X=x).
\]
If \(P(C\mid X=x)=0\), then both
\[
P(S_0=S_1,C\mid X=x)
\quad\text{and}\quad
P(C\mid X=x)
\]
are zero. If \(P(C\mid X=x)>0\), then \cref{obj:ass:weak-selection-monotonicity} gives one of the two almost-sure orders on the complier cell:
\[
1\{S_0=1\}\le 1\{S_1=1\}
\quad\text{or}\quad
1\{S_1=1\}\le 1\{S_0=1\}.
\]
The two indicators have equal conditional integrals by the preceding display, so the ordered indicators are equal almost surely on \(\{C,X=x\}\). Therefore
\[
P(S_0=S_1,\ C\mid X=x)=P(C\mid X=x).
\]
% lean: equalSelectionComplierMass_of_zero_gap

\item Under the same zero-gap condition, the argument in the first step also gives
\[
\Gamma_x^\ast=\Gamma_x^0.
\]
Indeed, every \(\gamma_x\in\Gamma_x^\ast\) has exact rows and exact columns because
\(q_0(x)=q_1(x)=m(x)\), and every \(\gamma_x\in\Gamma_x^0\) has bounded rows, bounded columns, and total mass \(m(x)\). Combining this equality with the already proved
\[
\Gamma_x^+=\Gamma_x^-,
\qquad
\Gamma_x^-=\Gamma_x^0,
\]
and with \(\Gamma_x=\Gamma_x^\ast\) from \cref{obj:def:partial-transport-polytope}, gives
\[
\Gamma_x=\Gamma_x^{+},
\qquad
\Gamma_x=\Gamma_x^{-},
\qquad
\Gamma_x=\Gamma_x^{0}.
\]
% lean: branchFree_eq_tie_of_gap_zero

\item Finally fix a threshold \(t\in\mathcal Y\). The prefix sums \(\ell_{\le t}(x)\) and \(h_{\le t}(x)\), and the upper tail \(h_{>t}(x)\), are those of \cref{obj:def:threshold-cuts}. The lower tail appearing in the statement is
\[
\ell_{>t}(x):=\sum_{\substack{i\in\mathcal Y\\ t<i}}\ell_i(x).
\]
Because \(\mathcal Y\) is finite and each index lies in exactly one of the two sets \(\{i:i\le t\}\) and \(\{i:t<i\}\),
\[
\ell_{\le t}(x)+\ell_{>t}(x)
=
\sum_{\substack{i\in\mathcal Y\\ i\le t}}\ell_i(x)
+
\sum_{\substack{i\in\mathcal Y\\ t<i}}\ell_i(x)
=
\sum_{i\in\mathcal Y}\ell_i(x)
=
q_0(x).
\]
The same finite partition gives
\[
h_{\le t}(x)+h_{>t}(x)
=
\sum_{\substack{j\in\mathcal Y\\ j\le t}}h_j(x)
+
\sum_{\substack{j\in\mathcal Y\\ t<j}}h_j(x)
=
\sum_{j\in\mathcal Y}h_j(x)
=
q_1(x).
\]
When \(\Delta q(x)=0\), \(q_0(x)=q_1(x)\), and hence
\[
\ell_{\le t}(x)-h_{\le t}(x)
=
\bigl(q_0(x)-\ell_{>t}(x)\bigr)
-
\bigl(q_1(x)-h_{>t}(x)\bigr)
=
h_{>t}(x)-\ell_{>t}(x).
\]
% lean: prefix_difference_eq_tail_difference_of_gap_zero
\end{enumerate}
The four displayed conclusions are exactly the asserted cellwise conclusions for every \(x\) with \(p_x>0\). % lean: tie_face_collapse
\end{proof}
